% ==========================================================================
%  Chapter 3: Rotations, Quaternions, and SO(3)
% ==========================================================================
\chapter{Rotations, Quaternions, and SO(3)}
\label{ch:rotations}

\epigraph{%
  The tragedy of Euclidean space is that it works perfectly until you spin.
}{}

\section{The Coordinate Frame}

When we model movement or control robotic systems, we are not looking at the physical object itself; we are tracking a \textbf{Coordinate Frame} rigidly attached to the object, often denoted as the Body Frame ($B$). 

To know where the Body Frame $B$ is relative to the absolute stationary universe—the Spatial (or World) Frame ($S$)—we must define two mathematical properties:
1. The **translation** vector $\bm{p} \in \Reals^3$ pointing from the origin of $S$ to the origin of $B$. 
2. The **rotation** of the axes of $B$ relative to the axes of $S$.

Unlike the straightforward $(X, Y, Z)$ translation vector, rotation is profoundly mathematically problematic.

\section{Euler Angles and Gimbal Lock}

The most intuitive way to describe the orientation of an aircraft, drone, or robotic arm is to split rotation into three successive angles around axes, usually denoted as Roll, Pitch, and Yaw ($X, Y, Z$). 

This is known as an \textbf{Euler Angle} representation: $\bm{\theta} = [\theta_1, \theta_2, \theta_3]^T$. While highly intuitive for a human pilot, Euler Angles harbor a fatal mathematical flaw known as \textbf{Gimbal Lock}.

\begin{laymansbox}{Gimbal Lock}{gimballock}
  If an airplane pitches perfectly $90^\circ$ straight up into the sky, its Roll axis (pointing out its nose) now perfectly aligns with its original Yaw axis (pointing perfectly down towards the Earth). 
  
  Mathematically, two independent degrees of freedom have become linearly dependent—they merged into the identical physical rotation axis. A degree of freedom is literally lost during the calculation. If you attempt to control the plane passing through this "Singularity", your matrices become non-invertible and your computer instantly crashes mapping a physical impossibility.
\end{laymansbox}

Since our control theory and geometric motion algorithms fundamentally rely on continuous trajectories twisting through all possible spatial angles (like chaotic golf club arcs or throwing motions), we must completely abandon Euler angles. We cannot afford mathematical singularities.

\section{The Special Orthogonal Group: SO(3)}

To map 3D orientation without singularities in all continuous space, we define a rotation by looking at the purely physical unit vectors defining the Body Frame $B$: the $\hat{x}_b, \hat{y}_b, \hat{z}_b$ axes.

If we write these three body axes explicitly in the coordinates of the Space Frame $S$, we can stack them vertically into a $3 \times 3$ matrix:
\begin{equation}
   R = \begin{bmatrix} | & | & | \\ \hat{x}_b & \hat{y}_b & \hat{z}_b \\ | & | & | \end{bmatrix} \in \Reals^{3 \times 3}
\end{equation}

This is a \textbf{Rotation Matrix}. 

Because $\hat{x}_b, \hat{y}_b,$ and $\hat{z}_b$ are all length $1$ (unit magnitude) and perfectly orthogonal (all $90^\circ$ apart), the columns are an orthonormal basis. The strict mathematical definitions that guarantee this orthonormality are:
1. $R^T R = I$ (The transpose equals the inverse)
2. $\det(R) = 1$ (The determinant is 1 to preserve right-handed orientation—no reflecting through a mirror)

\begin{principle}{The SO(3) Manifold}{so3}
  The continuous, infinite geometric space of all possible valid $3 \times 3$ rotation matrices satisfying $R^T R = I$ and $\det(R) = 1$ is a mathematical \textbf{manifold} called the \textbf{Special Orthogonal Group in 3 Dimensions}, denoted as $\SO$. 

  If a Configuration Space $\configspace$ is simply a rigid body rotating around a point, we say $\configspace = \SO$.
\end{principle}

\section{Unit Quaternions $\so$}

While $\SO$ completely prevents Gimbal Lock and provides pure, continuous rotation matrices, tracking 9 independent variables inside a $3 \times 3$ grid is wildly computationally inefficient for an engine updating a 15-DOF human model at 1,000 Hertz. 

To solve this, William Rowan Hamilton invented \textbf{Quaternions} ($\mathcal{H}$) in 1843. Rather than tracking a full Cartesian body frame $R \in \SO$, any 3D rotation can be flawlessly described as a combination of precisely two things:
1. A single unit vector $\hat{\bm{\omega}}$ defining the solitary physical axis of rotation.
2. An angle $\theta$ denoting how far to spin around that axis.

A Unit Quaternion $q$ mathematically encodes these into 4 variables representing an exponential scalar and vector map:
\begin{equation}
  q = \begin{bmatrix} q_0 \\ q_1 \\ q_2 \\ q_3 \end{bmatrix} = \begin{bmatrix} \cos(\theta/2) \\ \hat{\bm{\omega}}_x \sin(\theta/2) \\ \hat{\bm{\omega}}_y \sin(\theta/2) \\ \hat{\bm{\omega}}_z \sin(\theta/2) \end{bmatrix} \in \Reals^4
\end{equation}

By simply enforcing that the sum of the squares of these 4 variables equals 1 ($q_0^2 + q_1^2 + q_2^2 + q_3^2 = 1$), a Unit Quaternion provides a singularity-free, ultra-fast algebraic mapping of $\SO$.

Throughout *The Geometry of Motion*, specifically when interacting with the codebase bridges connecting algorithms like `Pinocchio` or `UpstreamDrift` (which natively uses MuJoCo physics engines), all 3D rotational mechanics explicitly bypass Euler Angles and map directly into Quaternions and $\SO$ Rotation Matrices.

To merge Translation and Rotation into a single unified control framework for physical bodies mapping the full world, we must upgrade from $3 \times 3$ Rotation Matrices to the magnificent $4 \times 4$ space of $\SE$.
